\section{Szenenüberblick}

\subsection{Einordnung in den Handlungsstrang}
Den Antagonisten in Fast X spielt \enquote{Dante Reyes}, verkörpert von Jason Momoa.
Als Sohn des früheren Drogenbarons \enquote{Hernan Reyes}\parencite{HernanReyesFFfandom} hat er dem Team um Dominic „Dom“ Toretto (gespielt von Vin Diesel) noch nicht verziehen. Dantes zentrale Motivation besteht in der Vergeltung für den Tod seines Vaters aus Fast and Furious 5 \parencite{HernanReyesFFfandom}. Diese Rachsucht prägt im übrigen die gesamte Filmhandlung 
\\
\\
%Rollenname in Anführungszeichen, Eigenname nicht; Nach Einführung des Spitznamen keine Anführungszeichen mehr nötig
Das Team wird zu Beginn des Films für eine dringende Mission nach Rom bestellt.
Dort sollen sie einen Supercomputer-Chip aus einem gepanzerten Konvoi extrahieren.
Als sie in dem LKW statt des Computerchips jedoch eine tonnenschwere \enquote{DM-79 sub-nautical neutron mine}\parencite{fastXwiki} vorfinden und der Lastzug plötzlich ferngesteuert wird, wird ihnen der Hinterhalt bewusst.
In einem spektakulären Versuch, das Fahrzeug am Eindringen in die Innenstadt zu hindern, überschlägt sich der Lastwagen, kommt zwar zum Halt, verliert dabei jedoch seine Ladung. 
\\
Die Bombe verlässt den LKW und rollt für die nächsten guten 10 Minuten ununterbrochen durch die Großstadt. Sie richtet massive Schäden an, und während Dom in seinem Auto sich aktiv darum bemüht, die Bombe durch physisches Eingreifen zu stoppen, prozediert Dante als heimlicher Drahtzieher, der die Bewegungen der Fahrzeuge orchestriert und den Fernzünder kontrolliert.
\\
Als Teil seines Racheplans nimmt er dabei bewusst den Tod von Doms Mitgliedern in Kauf.
Indem er aus dem Hintergrund agiert und die Helden ins Rampenlicht zwingt, verschleiert er die Verantwortung für diesen Anschlag geschickt. 
Durch die Inszenierung Doms als böswilligen Akteur lenkt er somit öffentlichen Unmut, und rechtliche Konsequenzen auf das Team\parencites{fastXfandom}{fastXwiki}.


%hochfrequente Actionszene versinnbildlichung film charakter 






\subsection{Drehort und Parameter}



\iffalse
Kugel rollt fast 8 Minuten Steigung in rom aber nur zwischen 120 und 13 über dem meerespiegel
https://en-gb.topographic-map.com/map-vbkn51/Rome/?center=41.88981%2C12.49526&zoom=12
https://www.imdb.com/title/tt5433140/goofs/?item=gf6927494&ref_=ext_shr_lnk

Kugel rollt an https://de.wikipedia.org/wiki/Torre_dei_Capocci und https://de.wikipedia.org/wiki/Spanische_Treppe

vorbei; sind mehere Kilometer auseinander
https://www.imdb.com/title/tt5433140/goofs/?item=gf8134727&ref_=ext_shr_lnk
\fi


Unschwer lässt sich erkennen, dass viele Teile dieser Szene in Rom gefilmt wurden, im Verlauf werden viele bekannte Sehenswürdigkeiten beschädigt und die Luftaufnahmen zeigen markante Hintergrunddetails. Gestützt wird diese Annahme durch die mehrere Aussagen bezüglich des Missionsziels, die im Film getätigt werden. Doch bei genauerer Untersuchung fallen einige Paradoxa auf. 
\\
Der Balkon, den das Team zu Beginn nutzt, um unbemerkt auf die Fahrzeuge zu springen und den Konvoi zu übernehmen, befindet sich nicht in Rom, sondern in der Via Vittorio Amedeo II 16 in Turin \parencite{TurinBalconyJump}. Die anschließende Verfolgungsjagd findet in einer Parallelstraße statt\parencite{FilmLocationsTurin}.
\\
Es wird kurz eine Luftaufnahme des Piazza del Popolo aus Rom\parencite{PiazzaDelPopoloWiki} gezeigt, bevor die Szene auf dem Piazza San Carlo in Turin mit dem Freisetzen der Bombe fortgesetzt wird \parencite{PiazzaSanCarloWiki}.
\\
Die folgenden Stationen finden dann tatsächlich in Rom statt, was aber nicht heißt, dass sich hier an die Gesetze der Raum-Zeit gehalten wird.
\\
Denn nachdem die Kugel unter Doms Einwirken nun fast vollständig zu Halt gekommen, und die Zwischenmission für einige Augenblicke geglückt scheint, wird sie in einer unglücklichen Abfolge von Ereignissen dann doch wieder in Bewegung gesetzt.  
Ein brennendes Auto, dass durch das Bild fliegt und per Kollision mit der Kugel für die Fortsetzung der Verfolgungsjagd sorgt, ist dabei zuerst vor dem Torre dei Cappocci \parencite{TorreDeiCapocciWiki} und im unmittelbar nächsten Schnitt vor der spanischen Treppe \parencite{SpanishStepsWiki} zu sehen. Diese Flugbahn ist komplett unmöglich da die beiden Sehenswürdigkeiten von mehreren Kilometern Innenstadt getrennt werden \parencite{imdbDistance}.
\\
Betrachtet man die Dauer der Rollsequenz im Verhältnis zur tatsächlich vorhandenen Höhendifferenz, offenbart sich ein deutlich gravierenderes physikalisches Paradoxon \parencites{imdbHeight}{topoHeightRome}.  
\\
Der vermeintliche Ausgangspunkt der Szene liegt in der Umgebung des Quirinal-Hügels \parencite{QuirinalWiki}, einem der höchsten Punkte Roms. Selbst unter Berücksichtigung dieser topographischen Gegebenheiten beträgt der Höhenunterschied zum Tiber jedoch kaum mehr als einige Dutzend Meter.  
\\
Eine Kugel, die unter idealisierten Bedingungen mehrere Minuten lang kontinuierlich beschleunigt, ist somit physikalisch kaum denkbar.  
\\
Vor diesem Hintergrund erscheinen die zuvor beschriebenen räumlichen Unstimmigkeiten fast vernachlässigbar, verglichen mit der Tatsache, dass Rom im Film offenbar über ein bislang unbekanntes, kilometerlanges Gefälle verfügt.
\\
\\
Zum Ende der Szene ist die Kugel fast bis in den Vatikan vorgedrungen. Um den Zerstörung des nahegelegenen Petersplatzes zu verhindern, versucht Dom nun die rollende Bombe zu überholen und sie durch vorsichtiges Bremsen direkt zu entschleunigen.
\\
Nachdem er sich des mangelnden Erfolgs dieser Strategie bewusst wird, ändert er jedoch seinen Plan. 
\\
Er beschleunigt spontan und steuert auf einen in unmittelbarer Nähe befindlichen Baukran zu. Er benutzt herumliegendes Baumaterial als provisorische Rampe und katapultiert sich dadurch in die Luft. Durch eine gezielte Kollision mit dem Ausleger wird der Kran in rotatorische Bewegung versetzt. Das Gegengewicht auf der anderen Seite des Drehpunkts schwingt in die Bahn der Kugel. Der folgende Zusammenstoß lenkt die Kugel von ihrer bisherigen Richtung ab und führt dazu, dass sie über die Brüstung der Brücke hinweg in den darunterliegenden Fluss stürzt. Kurz darauf detoniert sie, wodurch eine Druckwelle entsteht, deren Wirkung jedoch durch das Wasser erheblich abgeschwächt wird. Damit findet die Verfolgungssequenz ihr Ende.
\\
\\%TODO Screenshot behind the scenes; 
Auch dieser Abschnitt wurde jedoch nicht ausschließlich in Rom gedreht. Zwar verweist die gezeigte Szenerie eindeutig auf Rom, doch die Straßenführung und die Schnittkontinuität deuten an dieser Stelle klar auf Turin hin. Bei genauer Betrachtung lassen sich zudem Unstimmigkeiten in der Gestaltung der Brückenpfeiler erkennen.  
\\
Die \enquote{Behind the Scenes}-Aufnahmen erleichtern die Identifikation des tatsächlichen Drehorts. Im unbearbeiteten Rohmaterial ist deutlich zu sehen, dass die Brüstung, und damit auch die Brücke selbst, eine andere ist als in der finalen Filmfassung.  
\\
In der filmischen Realität befindet sich die Szene auf der Verbindungsbrücke Ponte Vittorio Emanuele II zwischen dem Zentrum Roms und dem Vatikan \parencite{RomePonteII}.  
Der tatsächliche Drehort war jedoch die Ponte Vittorio Emanuele I über dem Po in Turin \parencite{TurinPonteI}, die in der Nachbearbeitung digital angepasst wurde, um in weiten Teilen wie ihr namensverwandtes Pendant in Rom zu erscheinen.
\\
\\
Diese Widersprüche erschweren eine präzise Analyse der Umgebungsparameter, weshalb in den folgenden Teilen dieser Arbeit nur näherungsweise Rekonstruktionen vorgenommen werden können.


%Unkalibrierte Stereobild-Triangulation

\subsection{Ziel und mögliche Einflussnahme}


\subsection{Aufteilung in Untersequenzen}


\begin{wrapfigure}{r}{0.5\textwidth}
	\vspace{-10pt} % Optionale Anpassung des oberen Abstands
	\centering
	\includegraphics[width=0.48\textwidth]{resources/00_intro/sequences_blank_3x_nobg.png}
	%\vspace*{-10mm}
	\begin{flushright}
		{\scriptsize Eigenkreation; Gezeichnet in Excalidraw}
	\end{flushright}
	
	\captionsetup{
		labelfont=bf,
		textfont=footnotesize,
		skip=0pt
	}
	\caption{Aufteilung der Szene in abgeschlossene Untersequenzen}
	\label{fig:intro-sequences}
	\vspace{-10pt} % Optionale Anpassung des unteren Abstands
\end{wrapfigure}


\hyperref[fig:intro-sequences]{Abbildung~\ref*{fig:intro-sequences}} Zeigt die Unterteilung der Szenen 
\FloatBarrier